\SourceLeaf{219}%
\noindent Betrachtet man nun \(x\) und \(y\) als constant, und \(l,\lambda\) als veränderlich, so gibt die Differentiation
\[
\begin{aligned}
\frac{\mathrm{d}l}{\mathrm{d}\lambda}&=\frac{ab'-ba'}{(a\cos\lambda+a'\sin\lambda)^2+(b\cos\lambda+b'\sin\lambda)^2}\\
&=(ab'-ba')\mathord{.}\frac{\mathrm{d}\sigma^2}{\mathrm{d}s^2}
\end{aligned}
\]
\noindent Man sieht also, dass jenachdem \(ab'-ba'\) positiv \textit{oder} negativ ist, \(l\) und \(\lambda\) immer zugleich wachsen, \textit{oder} sich entgegengesetzt ändern, und also im erstern Fall die Darstellungen 2 und 3 ähnlich liegend, im andern verkehrt liegend sind.

Aus der Verbindung dieses Resultats mit dem vorhingefundenen ergibt sich, dass die Darstellungen in 1 und 3 ähnlich liegend oder verkehrt liegend sind, je nachdem \(\frac{ab'-ba'}{h}\) positiv oder negativ ist.

Da auf der Fläche, deren Gleichung \(\psi=0\) ist,
\[
e\,\mathrm{d}x+g\,\mathrm{d}y+h\,\mathrm{d}z=0
\]
\noindent also auch
\[
(ea+gb+hc)\mathrm{d}t+(ea'+gb'+hc')\mathrm{d}u=0
\]
\noindent wird, wie auch immer das Verhältniss von \(\mathrm{d}t\) und \(\mathrm{d}u\) gewählt wird, so muss offenbar identisch
\[
ea+gb+hc=0,\qquad ea'+gb'+hc'=0
\]
\noindent werden, woraus folgt, dass \(e,g,h\) resp. den Grössen \(bc'-cb',ca'-ac',ab'-ba'\) proportional sind, also
\[
\frac{bc'-cb'}{e}=\frac{ca'-ac'}{g}=\frac{ab'-ba'}{h}
\]
\noindent Man kann also, welchen dieser drei Ausdrücke man will, oder wenn man mit der ihrer Natur nach positiven Grösse \(ee+gg+hh\) multiplicirt, die sich ergebende symmetrische Grösse
\[
ebc'+gca'+hab'-ecb'-gac'-hba'
\]
\noindent als Criterium der ähnlichen oder verkehrten Lage der Theile in den Darstellungen 1 und 3 anwenden.

Ganz eben so wird ähnliche oder verkehrte Lage der Theile in den Darstellungen 6 und 8 von dem positiven oder negativen Werthe der Grösse
\[
\frac{BC'-CB'}{E}=\frac{CA'-AC'}{G}=\frac{AB'-BA'}{H}
\]
